\documentclass[12pt]{article}
\everymath{\displaystyle}
\usepackage{unicode-math}
\usepackage{pgfplots}
\pgfplotsset{compat=1.18}
\usepackage{amsmath}



\usepackage[margin=1in]{geometry}
\usepackage{tcolorbox}
\usepackage{graphicx}
\usepackage{tikz}
\usepackage{amssymb}
\usepackage{enumitem}
\usepackage{titlesec}
\usepackage{fancyhdr}
\usepackage{pifont}
\usepackage{xcolor}
\newcounter{questionnum}
\setcounter{questionnum}{0}
\usepackage{eso-pic} % for page background

% --------- 主题颜色设置 ---------
\definecolor{novelblue}{RGB}{70,60,150}
\definecolor{headerbg}{RGB}{240,240,255}

\pagestyle{fancy}
\fancyhf{}
\renewcommand{\headrulewidth}{0.4pt}
\renewcommand{\headrule}{\hbox to\headwidth{\color{novelblue}\leaders\hrule height 0.4pt\hfill}}

% --------- 页眉背景色条 ---------
\newcommand{\headbackground}{
  \AddToShipoutPictureBG*{
    \begin{tikzpicture}[remember picture, overlay]
      \fill[blue!5!purple!10] (current page.north west) rectangle ([yshift=-60pt]current page.north east);
    \end{tikzpicture}
  }
}

% --------- 页眉设置 ---------
\pagestyle{fancy}
\fancyhf{}
\renewcommand{\headrulewidth}{0pt}
\fancyhead[L]{\color{novelblue}\textbf{\textbf{Novel Prep} }}
\fancyhead[C]{\color{novelblue}\textbf {SAT Preparation}}
\fancyhead[R]{\color{novelblue}\textbf{Jack Zeng}}


\newtcolorbox{SATCard}[1][]{%
  enhanced,
  colback=white,
  colframe=novelblue,
  coltitle=white,
  boxrule=0.6pt,
  arc=4pt,
  boxsep=5pt,
  left=8pt, right=8pt, top=10pt, bottom=10pt,
  sharp corners=south,
  drop shadow south east,
  #1
}

% --------- 顶部 Mark for Review 栏 ---------
\newcommand{\markforreview}{
  \stepcounter{questionnum} % 自动递增题号
  \begin{tikzpicture}[scale=1.1]
    % 黑色题号
    \fill[black] (0,0) rectangle (0.8,0.6);
    \node[white, font=\bfseries\large] at (0.4,0.3) {\thequestionnum};

    % 灰底主栏
    \fill[gray!10] (0.8,0) rectangle (14.5,0.6);

    % Mark for Review
    \node[anchor=west, font=\small] at (1.2,0.3) {\ding{72} \hspace{2pt}Mark for Review};

    % ABC 按钮区域
    \draw[thick, rounded corners=3pt] (13.5,0.12) rectangle (14.5,0.48);
    \node[font=\bfseries\normalsize] at (14.0,0.3) {ABC};
  \end{tikzpicture}


  \newcommand{\choiceboxcircled}[2]{%
  \begin{tcolorbox}[
    colback=white,
    colframe=black,
    boxrule=0.5pt,
    rounded corners=6pt,
    left=6pt, right=6pt, top=6pt, bottom=6pt,
    width=\linewidth
  ]
    \textbf{\large\textcircled{\textbf{#1}}} \ #2
  \end{tcolorbox}
}
}

% --------- 选项框样式 ---------
\newcommand{\choicebox}[2]{%
  \begin{tcolorbox}[
    colback=white, 
    colframe=novelblue,
    boxrule=0.5pt, 
    rounded corners=4pt, 
    left=6pt, right=6pt, top=6pt, bottom=6pt,
    width=\linewidth
  ]
  \textbf{#1.} \ #2
  \end{tcolorbox}
}



\begin{document}
\section*{SAT Hard Question 17}


\noindent\markforreview
\vspace{1em}

For the function $f$, for each increase in the value of $x$ by $c$, where $c$ is a positive constant, the value of $f(x)$ increases by a factor of $27$. Which of the following equivalent forms of the function $f$ displays $\frac{1}{c}$ as a coefficient of $x$?

\begin{enumerate}[label=(\Alph*)]
\item $f(x) = 54(3)^{\frac{1}{3}x}$
\item $f(x) = 54(3^3)^{\frac{1}{6}x}$
\item $f(x) = 54(9)^{\frac{1}{4}x}$
\item $f(x) = 54\left(27^{\frac{1}{3}x} \right)^{\frac{1}{2}}$
\end{enumerate}



\vspace{2em}
\noindent\markforreview
\vspace{1em}

The composition of an animal is defined as the muscles, bones, and fat of the animal. A scientist studied the composition of one young swamp buffalo and determined the buffalo had $133$ kilograms of muscle, which made up approximately $65.9\%$ of its composition. Of the remaining composition of this buffalo, approximately $46.2\%$ was bone, and the remainder was fat. Based on these approximations, to the nearest tenth, how many kilograms of this buffalo's composition was bone?
\newpage
\vspace{2em}
\noindent\markforreview


\[f(x) = 3(x-a)(x-b)(x-c)\]


The function $f$ is defined by the given equation, where $a$, $b$, and $c$ are distinct constants. When $a < x < b$, the value of $f(x)$ is positive. The graph of $y = f(x)$ in the $xy$-plane contains the point $(r, s)$, where $r$ and $s$ are constants. If $s = 8$, which of the following could be true?

\begin{itemize}
\item[I.] $r < a$
\item[II.] $b < r < c$
\item[III.] $r > c$
\end{itemize}

\begin{enumerate}[label=(\Alph*)]
\item I only
\item III only
\item I and III only
\item II and III only
\end{enumerate}


\vspace{2em}
\noindent\markforreview
\vspace{1em}

In isosceles triangle $ABC$, sides $AB$ and $AC$ are congruent. Point $D$ divides side $BC$ such that the length of $\overline{BD}$ is $\dfrac{4}{7}$ of the length of $\overline{BC}$. Point $E$ lies on side $AB$ and point $F$ lies on side $AC$ such that when segments $\overline{DE}$ and $\overline{DF}$ are drawn, angle $\angle BED$ is congruent to angle $\angle CFD$. If the length of $\overline{BE}$ is $8$, what is the length of $\overline{CF}$?

\begin{itemize}
  \item[(A)] 6
  \item[(B)] 14
  \item[(C)] 32
  \item[(D)] 56
\end{itemize}

\newpage
\vspace{2em}
\noindent\markforreview
\vspace{1em}

The function $f$ is defined by the given equation: 
\[
f(x) = x^2 - 4x - 320
\]
Which of the following equivalent forms of the equation displays the minimum value of the function as a constant or coefficient?

\begin{itemize}
  \item[(A)] $f(x) = x^2 - 4(x + 80)$
  \item[(B)] $f(x) = (x - 2)^2 + (-324)$
  \item[(C)] $f(x) = x(x - 4) + (-320)$
  \item[(D)] $f(x) = (x + 16)(x - 20)$
\end{itemize}



\vspace{2em}
\noindent\markforreview
\vspace{1em}

In the given equation, $s$ and $r$ are constants, and $s > 0$. If the equation has infinitely many solutions, what is the value of $s$?
\[
\frac{12x + 28}{4} - \frac{s}{13} = r(x - 8)
\]
\vspace{2em}


\vspace{2em}
\noindent\markforreview
\vspace{1em}

A right triangular pyramid has exactly six edges. Each of these edges is $96$ centimeters long. If the surface area of the pyramid is $k\sqrt{3}$ square centimeters, what is the value of $k$?
\vspace{2em}







\end{document}